\documentclass[a4paper]{article}
\usepackage{amsmath}
\usepackage{amssymb}
\usepackage{amsthm}
\usepackage{geometry}
\usepackage{hyperref}

\newcommand{\Jnorm}{J_{\mathrm{norm}}}
\newtheorem{lemma}{Lemma}
\newtheorem{proposition}[lemma]{Proposition}
\newtheorem{definition}[lemma]{Definition}

\geometry{margin=25mm}

\title{Neutron \(\beta\) Decay Is a Replacement of Rotor Class:\\
The Quark-Charge Chart}

\author{https://github.com/hypernumbernet}
\date{\today}

\begin{document}
\maketitle

\begin{abstract}
Neutron \(\beta\) decay is read as a replacement of the two nucleon classes of double spacetime theory. The discrete companion writes a nucleus as an ordered product of the proton class \([R_{\rm p}]\) and the neutron class \([R_{\rm n}]\), and assigns each class baryon number \(B=1\) and dual-phase winding \(w=3\). The strong companion takes the neutron to be the case \(e=0\) of the electromagnetic label and the proton to be the case in which that label is on. Free-neutron decay replaces one factor of the product, from \([R_{\rm n}]\) to \([R_{\rm p}]\), keeps the number of factors and the winding, and raises the electromagnetic label of that factor to \(+e\). The emission that preserves the total charge is the electron of charge \(-e\) already used by the shell companion. Continuity of the electron energy spectrum is the laboratory requirement of a third, electrically silent body. The six-generator space has room for that body. The replacement itself does not force the body, because both nucleon classes can sit on the balanced locus \(J=0\).

The quark charges \(+2/3\) and \(-1/3\) are the unique rational solution that reproduces the two nucleon charges \(+1\) and \(0\) with three additive labels and only two values. The chart is a coordinate that apportions the integer charges of the two classes onto three labels. The three spatial axes of the parent theory are three embeddings of one and the same \(3+3\) split. Projectors of distinct axes do not commute and do not sum to the unit, so the axes are neither three simultaneous colors nor three generations. Compact \(\mathrm{SU}(2)\) is the dual rotor. The color-counting factor is not read off from this split.

The free mass gap \(m_n-m_p\) lies strictly between \(1293/1000\,\mathrm{MeV}\) and \(1294/1000\,\mathrm{MeV}\). The first node of the strong companion lies strictly between \(14/19\,\mathrm{fm}\) and \(213/280\,\mathrm{fm}\), and \(ke^{2}/r\) at that radius is strictly greater than \(189/100\,\mathrm{MeV}\). A positive electric self-energy makes the charged class heavier. The free ordering is the other way, and the gap is smaller than that self-energy. The class ordering is not the self-energy at the node. The mean life lies between \(870\,\mathrm{s}\) and \(890\,\mathrm{s}\), while \(\hbar/(m_n-m_p)c^{2}\) is less than \(6\times 10^{-22}\,\mathrm{s}\). The ratio exceeds \(10^{24}\). The replacement is not a crossing of a node of \(\gamma_s\).

The V--A coupling, the lifetime, the shape of the spectrum, \(g_A/g_V\), the Fermi constant, the neutrino mass, the generations, the Cabibbo--Kobayashi--Maskawa matrix, the structure functions, the ratio \(R\), and the absolute mass gap are not derived.
\end{abstract}

\section{Introduction}
\label{sec:intro}

The parent theory reads the four interactions as two channels of one geometry~\cite{dst-main}. The gravitational channel is the pair \((\phi,\theta)\) tied to inertial mass, and the electromagnetic channel is the pair \((\alpha,\beta)\) tied to charge. The strong force is proposed as the layered phase of the gravitational channel at nuclear density. Color charge is named there as a misreading of the torsional charge carried by dual-rotor phases, and the coupling \(\alpha_s\) is not identified with the discrete floor \(\varepsilon_N\). Electromagnetism and the weak channel occupy the same six-generator space, and neither the dual map nor the dual rotor supplies a V--A coupling.

The discrete companion embeds each nucleon as a three-layer torsion star~\cite{dst-dio}. The standard classes \([R_{\rm p}]\) and \([R_{\rm n}]\) carry \(B=1\) and \(w=3\), and a nucleus of mass number \(A\) is an ordered product of their representatives. Algebraic stability requires the product to lie in the unit group, the winding to satisfy \(w=3\ell\) for an integer \(\ell\) compatible with the three-layer topology, and \(|\Jnorm|\le 1\). The strong companion uses the same two classes as the particles of the nuclear force~\cite{dst-strong}. The neutron carries only the gravitational channel, and the proton adds the electromagnetic term. A proton and a neutron are different classes, so a proton--neutron pair at one separation is not forbidden by the exclusion of a shared seat of identical fermions.

This note reads the transition between those two classes as neutron \(\beta\) decay. Section~\ref{sec:chart} shows that the quark charges are a chart of the two integer charges. Section~\ref{sec:axes} separates the three axes from the three particles of that chart. Section~\ref{sec:classes} collects the quantities the two classes already fix. Section~\ref{sec:decay} is the list of what the replacement conserves. Section~\ref{sec:gap} shows that the free mass gap is not the Coulomb energy at the node. Section~\ref{sec:moment} records the ratios the same chart assigns to the magnetic moments, and the gaps by which the measured ratios miss them. Section~\ref{sec:open} lists what this note does not derive about the weak channel.

\section{The quark-charge chart}
\label{sec:chart}

The laboratory charges of the two nucleons are \(+e\) for the proton and \(0\) for the neutron. The quark model writes these two as sums of three additive charges and assumes that the charges take only two values, \(u\) and \(d\). Assigning two and one to the proton, and one and two to the neutron, gives the chart
\begin{equation}
\label{eq:chart}
2u+d=1,
\qquad
u+2d=0,
\end{equation}
in units of \(e\).

\begin{lemma}
\label{lem:chart}
The only pair of rationals \((u,d)\) satisfying~\eqref{eq:chart} is \(u=2/3\), \(d=-1/3\).
\end{lemma}

\begin{proof}
The difference of the two equations is \(u-d=1\). The second equation gives \(u=-2d\), so \(-2d-d=1\), hence \(d=-1/3\) and \(u=2/3\).
\end{proof}

The pair does not make every three-label sum an integer. The sum is \(+1\) or \(0\) only for the assignments \((2,1)\) and \((1,2)\). Three labels \(u\) give \(+2\), and three labels \(d\) give \(-1\). The chart keeps the charges of free particles integral because it restricts nucleons to those two assignments. The meson chart adds one value and its opposite and is likewise integral: the sum of \(u\) and \(\bar d\) is \(+1\), and the sum of \(u\) and \(\bar u\) is \(0\). The fractions cancel inside the restricted assignments.

The lemma says that the fractional charges are the unique inversion of the two nucleon charges. Once three additive carriers are assumed, the integers measured in the laboratory force \(2/3\) and \(-1/3\). Without that assumption, what is forced remains the integers. A nucleon of double spacetime theory is one rotor class, not a product of three carriers. Section~\ref{sec:axes} gives the reason the three axes cannot be read as those carriers. The chart remains as a coordinate in which the charges of the two classes are written with three labels. Relations that propagate the labels as particles --- structure functions, jets, and the ratio \(R\) --- lie outside the lemma.

The chart value of \(R\), with a color count put in separately, is \(R=N_c\sum_f q_f^2\). For \(u\), \(d\), \(s\) and \(N_c=3\) one has \(R=2\). Without the count the sum is \(2/3\). This note does not recompute the plateaus. The counting factor cannot be read off from the three axes either, by the obstruction of the next section.

\section{Three axes}
\label{sec:axes}

Relative to one spatial axis \(e_a\), the parent theory splits the six generators into a commuting triple \((B_a,E_{a+1},E_{a+2})\) and an anticommuting triple \((E_a,B_{a+1},B_{a+2})\)~\cite{dst-main}. The commuting triple is \(\mathfrak{so}(2,1)\), not \(\mathfrak{su}(2)\). The only triple in which every element squares to \(-1\) is the magnetic triple, and those three satisfy the quaternion table \(IJ=K\). Compact \(\mathrm{SU}(2)\) in this six-space is the dual rotor.

Each of the three spatial axes furnishes such a split. The associated projectors are idempotent. Projectors of distinct axes do not commute, and the three right projectors do not sum to the unit. The parent theory states them as alternative embeddings of one and the same \(3+3\) geometry, and not as three generation labels. The same sentence closes color. Three simultaneous colors are labels that commute with one another and exhaust the unit inside one representation. The axis projectors have neither property. The three axes are therefore not three quarks, and they are not three colors.

The winding \(w=3\) stands on the other side of this obstruction. What the discrete companion assigns to a nucleon is not baryon number \(1/3\) for each of three particles, but the closing of the dual phase of one class in three layers. A transition that sent one axis outward would not change the number of factors in the ordered product. An axis is not a factor. The asymptotic states counted by the discrete theory are the classes \([R_{\rm p}]\) and \([R_{\rm n}]\), not axes. The fractions of the chart do not appear as free particles because, in this sense, there is no particle in the product that would have to be confined. This is not a derivation of a confining potential. The exponential screening that the strong companion states inside the node is a separate statement about the force on a nucleon, and this note does not identify that screening with asymptotic freedom.

The same obstruction already stands in front of the generations. The parent theory does not read a finite discrete rotor image as a count of generations. The muon and the tau are not the other two axes of the nucleon. The branching by which a charged pion yields a muon more often than an electron is a fact about generations, and it is not fixed by the replacement of this note.

\section{The two nucleon classes}
\label{sec:classes}

The properties of the classes used here are only those already fixed by the companions.

\begin{enumerate}
\item Both carry baryon number \(B=1\) and winding \(w=3\)~\cite{dst-dio}. One nucleon closes in three layers, and the integer \(\ell\) in \(w=3\ell\) is \(\ell=1\) for a nucleon.
\item The electromagnetic label of the neutron is \(e=0\), and the electromagnetic label of the proton is \(+e\)~\cite{dst-strong}. Charge is a label of a channel, not the sign of the mismatch \(J\).
\item On the reading of nuclear matter, the equal-scale locus has \(J=0\) and \(M>0\)~\cite{dst-main}. A balanced seed is not emptied by either rapidity budget alone.
\item The electron is a particle of the electromagnetic channel, and its rest is held by the equality of the two budgets~\cite{dst-shell}. A hydrogen ionisation does not pay that rest.
\end{enumerate}

What is not yet fixed as a difference of classes is the difference of unsigned mass \(M\). That the free neutron is heavier than the free proton is a measurement, and the next section uses it as an input. No mass is defined per axis. The excess that the chart writes as \(m_d>m_u\) is, here, an excess of the class \([R_{\rm n}]\), not an excess of the label \(d\).

The ordered product has no third nucleon class. A decay that changes strangeness requires a class this note does not introduce. Hyperons lie outside it.

\section{Beta decay as replacement of class}
\label{sec:decay}

\begin{definition}
\label{def:replacement}
A \emph{replacement of class} in the ordered product of a nucleus is the replacement of one factor \([R_{\rm n}]\) by one factor \([R_{\rm p}]\), or the reverse.
\end{definition}

\begin{proposition}
\label{prop:conserved}
A replacement of class preserves the number \(A\) of factors and preserves the winding condition \(w=3\) of each factor. The replacement of a neutron by a proton changes the electromagnetic label of that factor from \(0\) to \(+e\). The reverse returns the label from \(+e\) to \(0\).
\end{proposition}

\begin{proof}
By the definition of the ordered product, \(A\) is the number of factors. A replacement exchanges one factor for one factor, so the number is unchanged. Both classes carry \(w=3\), so the closed-torsion-loop condition holds after the replacement if it held before. The electromagnetic labels are those the strong companion assigns to the two classes.
\end{proof}

Read as the number of factors, baryon number is preserved by the replacement. There is no need to change one label of baryon number \(1/3\). Because the winding is preserved, the decay is not an unwinding of one of the three layers.

If the total charge is conserved, the replacement of a neutron by a proton sends charge \(-e\) out of the product. That charge is the label already carried by the electron of the shell companion. The electron is not a fragment of the nucleon. It is the electromagnetic rotor that leaves in order to keep the total when the replacement turns the nucleon's electromagnetic label on.

The reverse replacement has two laboratory realizations. In one, the product emits a positron. In the other, it absorbs an electron that is already present. The first is \(\beta^{+}\) decay, and the second is electron capture. Both are the reverse of Proposition~\ref{prop:conserved}, and both return the electromagnetic label to silence. The branching ratio is not derived.

The third body does not follow from the conservation laws of the replacement. Continuity of the electron kinetic-energy spectrum is the laboratory requirement that an electrically silent particle leaves with the electron. The six-generator space has room for a dual-sector excitation with no charge label, as a reading of that particle. The forcing is a separate matter.

\begin{proposition}
\label{prop:no-third}
When both nucleon classes can sit at \(J=0\), the mismatch carried out by the replacement is zero. A charged rotor can sit at \(J=0\), \(M>0\). A balanced charged emission therefore does not require a second rotor to cancel the mismatch.
\end{proposition}

\begin{proof}
Nuclear matter in the parent theory has \(J=0\) on the equal-scale locus. If both classes lie on that locus, replacement of class does not change the nucleon's \(J\). The same theory distinguishes a configuration with \(E^{2}=B^{2}\), which is a balanced mass \(J=0\), \(M>0\), from a pure electric configuration. The rest of the electron is read as that configuration, the one in which the two budgets agree. The charge label is independent of the sign of \(J\). Charge \(-e\) and \(J=0\) are therefore compatible on one rotor, and a magnetic rotor assigned to cancel the mismatch is not required by the balance condition.
\end{proof}

The silent particle required by the continuous spectrum remains a laboratory fact after this proposition. What the proposition removes is the reading that derives the particle from the balance of \(J\). Assigning the particle a lepton number, identifying it with the neutrino and then fixing a mass, and fixing the chiral structure of the emission all lie outside this six-space. The parent theory has already shown that the dual map is not a chirality projector, and that the parallel plane contains no chirality-preserving Weinberg angle. This note does not use those two negations as a substitute for V--A.

In free space the direction of the replacement is fixed by the masses. The rest energy of the free neutron is greater than the sum of the rest energies of the free proton and the electron. The excess is the gap of Section~\ref{sec:gap} with the electron rest subtracted, and it lies strictly between \(782/1000\,\mathrm{MeV}\) and \(783/1000\,\mathrm{MeV}\). The free-neutron replacement is therefore open, and the reverse replacement of a free proton is closed. This opening and closing follows from the masses as inputs. The algebra of the classes does not choose which one is heavier.

Inside a nucleus the Coulomb sum of the strong companion can reverse the direction. In every chamber with \(\gamma_E>0\), adding one proton adds a Coulomb sum over the proton pairs already present, and adding a neutron does not increase that sum. When the electromagnetic cost of a proton-rich product exceeds the free gap, the reverse replacement lowers the energy of the nucleus. This note does not evaluate the nuclide at which that crossing occurs. The coefficient of \(Z^{2}/A^{1/3}\) in the semi-empirical mass formula remains the coefficient the strong companion already does not derive.

\section{The sign of the mass gap}
\label{sec:gap}

The CODATA free masses place the difference \(m_n-m_p\) between \(12933/10000\,\mathrm{MeV}\) and \(12934/10000\,\mathrm{MeV}\), and the electron rest energy between \(51099/100000\,\mathrm{MeV}\) and \(51101/100000\,\mathrm{MeV}\)~\cite{codata2018}. The ratio \((m_n-m_p)/m_e\) therefore lies strictly between \(5/2\) and \(38/15\). The lower end is the smaller gap divided by the larger electron mass, \(1.2933/0.51101>5/2\), and the upper end is the larger gap divided by the smaller electron mass, \(1.2934/0.51099<38/15\). The interval does not identify the ratio.

\begin{proposition}
\label{prop:sign}
In the range of lengths of the strong companion, \(ke^{2}/r\) at the first node is strictly greater than \(189/100\,\mathrm{MeV}\). The free mass gap is strictly smaller, and the sign is the other way.
\end{proposition}

\begin{proof}
From the laboratory signs, the companion places the first node strictly between \(14/19\,\mathrm{fm}\) and \(213/280\,\mathrm{fm}\). The quantity \(ke^{2}=\alpha\hbar c\) is smallest at the upper end of this interval. One has \(\alpha\hbar c>1.439\,\mathrm{MeV\,fm}\), and
\[
1.439\times\frac{280}{213}>1.891>1.89.
\]
The mass gap is less than \(1.294\,\mathrm{MeV}\). A positive \(ke^{2}/r\) raises the energy of the charged side. In the free ordering the charged class is the lighter one.
\end{proof}

The excess of the class \([R_{\rm n}]\) is therefore not the value of the proton's electric self-energy placed at the node. The self-energy works on the side that would cancel the observed ordering, and it is larger than the gap. The excess belongs to the unsigned mass of the class itself, and not to the Coulomb term. The absolute value of that mass is not derived. The chart's \(m_d-m_u\) is the same excess apportioned onto a label, and it is not derived here either.

The time scale likewise separates the replacement from the strong channel. If \(\hbar c<198\,\mathrm{MeV\,fm}\) and \(c>2.99\times 10^{23}\,\mathrm{fm\,s^{-1}}\), then \(\hbar<6.63\times 10^{-22}\,\mathrm{MeV\,s}\). Since the gap is greater than \(1.293\,\mathrm{MeV}\),
\[
\frac{\hbar}{(m_n-m_p)c^{2}}<\frac{6.63\times 10^{-22}}{1.293}\,\mathrm{s}<6\times 10^{-22}\,\mathrm{s}.
\]
The reported mean life lies between \(870\,\mathrm{s}\) and \(890\,\mathrm{s}\)~\cite{pdg2024}. The ratio exceeds \(10^{24}\). The nuclear times of order MeV treated by the strong companion do not lie on the side of this lifetime. Beta decay is not an infall across a node of \(\gamma_s\).

\section{The magnetic-moment chart}
\label{sec:moment}

The same chart fixes ratios of magnetic moments once spin weights are assigned to the three labels. With the standard weights,
\begin{equation}
\label{eq:moments}
\mu_p=\frac43\mu_u-\frac13\mu_d,
\qquad
\mu_n=\frac43\mu_d-\frac13\mu_u,
\end{equation}
and the charge ratio together with equal masses gives \(\mu_u=-2\mu_d\).

\begin{lemma}
\label{lem:moments}
If \(\mu_u=-2\mu_d\) and \(\mu_d\neq 0\), then~\eqref{eq:moments} yields \(\mu_n/\mu_p=-2/3\) and \((\mu_p+\mu_n)/(\mu_p-\mu_n)=1/5\).
\end{lemma}

\begin{proof}
Substitution gives \(\mu_p=(3/2)\mu_u\) and \(\mu_n=-\mu_u\). The ratio is \(-2/3\). The sum is \(\mu_u+\mu_d=-\mu_d\), and the difference is \((5/3)(\mu_u-\mu_d)=(5/3)(-3\mu_d)=-5\mu_d\). The ratio of sum to difference is \(1/5\).
\end{proof}

The measurements lie near these two rationals and not on them. In nuclear magnetons, CODATA places \(\mu_p\) between \(2.79284\) and \(2.79286\), and \(\mu_n\) between \(-1.91305\) and \(-1.91303\)~\cite{codata2018}. On this rectangle the ratio \(\mu_n/\mu_p\) lies strictly between \(-1.91305/2.79284\) and \(-1.91303/2.79286\), and therefore strictly between \(-137/200\) and \(-171/250\). Its difference from \(-2/3\) lies strictly between \(1/55\) and \(1/54\). The sum \(\mu_p+\mu_n\) lies between \(0.87979\) and \(0.87983\), the difference \(\mu_p-\mu_n\) lies between \(4.70587\) and \(4.70591\), and their ratio lies strictly between \(9/50\) and \(19/100\). The deficit from \(1/5\) lies strictly between \(1/80\) and \(1/70\).

The weights \(4/3\) and \(-1/3\) do not come from the \(\mathfrak{so}(2,1)\) stabilizer. The only compact generator in the stabilizer is \(B_a\), which holds the axis, and it does not classify three carriers of spin \(1/2\). This note does not improve the ratio \(-2/3\). It records both that the chart comes close and that the geometry of the axes does not supply the ratio.

\section{What is not derived}
\label{sec:open}

The following quantities follow neither from the replacement of class nor from the lemmas of the chart.

The lifetime of the weak decay, and the shape of the electron spectrum. The Fermi constant remains an input. The factor \(10^{24}\) above is a comparison showing that the lifetime is not the inverse time of the mass gap. It is not a value of \(G_F\).

The ratio \(g_A/g_V\). The axial coupling is not read off from the sandwich in the six-space. What the sandwich supplies is the Lorentz four-force of \((E,-B)\), which the parent theory already distinguishes from V--A.

The neutrino mass, its helicity, and lepton number. Proposition~\ref{prop:no-third} does not derive the silent particle from balance.

The three generations, the Cabibbo angle, and the Kobayashi--Maskawa matrix. The three axes are not generations. The ordered product has only two classes and contains no flavor-changing factor.

Structure functions and the ratio \(R\). Lemma~\ref{lem:chart} is an inversion of the integer charges, not a momentum fraction delivered by a virtual photon to one axis. The color count cannot be read as the count of axes.

The absolute mass gap, and \(\alpha_s\). Proposition~\ref{prop:sign} is an obstruction of sign and of scale. It is not a derivation of \(1.293\,\mathrm{MeV}\).

Charged-pion decay stands in the same channel as the transition in which a winding-zero object emits its electromagnetic label as a lepton and leaves no nucleon factor. The strong companion already uses the charged-pion Compton wavelength as an identification of the length \(\ell\), and that identification is not a derivation. That the pion selects the muon over the electron does not follow from here either.

\section{Conclusion}
\label{sec:conclusion}

Neutron \(\beta\) decay is the transition that replaces one factor of the ordered product, from \([R_{\rm n}]\) to \([R_{\rm p}]\). The number of factors and the winding \(w=3\) are preserved, and the electromagnetic label of that factor changes from \(0\) to \(+e\). The emission that preserves charge is the electron. The electrically silent particle required by the continuous spectrum has a place to be read in the six-space, and it is not forced by the balance \(J=0\). In free space the replacement runs from neutron to proton because the neutron class is heavier than the sum of the proton class and the electron, and that inequality is an input. The reverse can open inside a nucleus because the Coulomb sum over proton pairs can exceed the free gap.

The quark charges \(2/3\) and \(-1/3\) are the unique rational solution left when the two nucleon charges are written with three additive labels and two values. The three axes are not the three particles of those labels. They are alternative embeddings of one \(3+3\) split, and they are neither colors nor generations. The winding \(3\) is the closing of one class in three layers. The free mass gap is smaller than the positive Coulomb energy placed at the first node, and the sign is the other way. The magnetic-moment chart yields \(-2/3\) and \(1/5\), and the measurements leave gaps on the smaller side of both, strictly between \(1/55\) and \(1/54\), and strictly between \(1/80\) and \(1/70\). The lifetime is longer than the inverse time of the mass gap by more than \(10^{24}\).

The operator of the weak interaction remains underived after this reading.

\begin{thebibliography}{9}

\bibitem{dst-main}
Anonymous,
``Gravity as Torsion between Dual Spacetime: A Biquaternionic Reformulation of General Relativity'' (2025).
\url{https://github.com/hypernumbernet/dual-spacetime-doc/blob/main/double-spacetime-theory.pdf}.

\bibitem{dst-dio}
Anonymous,
``Discrete Biquaternionic Double Spacetime: Rigorous Proofs of Diophantine Conjectures via Torsional Mismatch and PGA Embedding,''
companion note.

\bibitem{dst-strong}
Anonymous,
``The Strong Force as a Repulsive Layer of Proximity Gravity,''
companion note.

\bibitem{dst-shell}
Anonymous,
``Electron Orbits from the Signed Electrostatic Force: Covalent Bonds, Metallic Bonds, and Crystals,''
companion note.

\bibitem{codata2018}
E.~Tiesinga, P.~J.~Mohr, D.~B.~Newell, and B.~N.~Taylor,
``CODATA recommended values of the fundamental physical constants: 2018,''
\textit{Rev.\ Mod.\ Phys.} \textbf{93}, 025010 (2021).

\bibitem{pdg2024}
S.~Navas \textit{et al.} (Particle Data Group),
``Review of Particle Physics,''
\textit{Phys.\ Rev.\ D} \textbf{110}, 030001 (2024).

\end{thebibliography}

\end{document}
